% Compile from delivery/_src:
% tectonic --outdir .. handwritten-answers-tex.tex
\documentclass[11pt,a4paper]{article}
\usepackage[margin=18mm,headheight=15pt]{geometry}
\usepackage{fontspec}
\setmainfont{Times New Roman}
\setsansfont{Arial}
\setmonofont{Courier New}[Scale=0.9]
\newfontfamily\cyrillicfont{Times New Roman}[Script=Cyrillic]
\newfontfamily\cyrillicfontsf{Arial}[Script=Cyrillic]
\newfontfamily\cyrillicfonttt{Courier New}[Script=Cyrillic,Scale=0.9]
\usepackage{polyglossia}
\setdefaultlanguage{russian}
\setotherlanguage{english}
\usepackage{amsmath,amssymb,mathtools}
\usepackage{booktabs,tabularx,array}
\usepackage{enumitem}
\usepackage{xcolor}
\definecolor{ink}{HTML}{1D3557}
\definecolor{muted}{HTML}{666666}
\usepackage{titlesec}
\titleformat{\section}{\Large\sffamily\bfseries\color{ink}}{}{0pt}{}
\titleformat{\subsection}{\normalsize\sffamily\bfseries\color{ink}}{}{0pt}{}
\titlespacing*{\section}{0pt}{0pt}{10pt}
\titlespacing*{\subsection}{0pt}{11pt}{5pt}
\usepackage{fancyhdr}
\pagestyle{fancy}
\fancyhf{}
\fancyhead[L]{\small\sffamily Machine Learning -- Semester 2}
\fancyhead[R]{\small\sffamily Конспекты по фотографиям}
\fancyfoot[C]{\small\thepage}
\renewcommand{\headrulewidth}{0.3pt}
\usepackage[unicode,hidelinks,pdfusetitle]{hyperref}
\hypersetup{pdftitle={Machine Learning - Semester 2: typeset handwritten notes}}
\setlength{\parindent}{0pt}
\setlength{\parskip}{5pt}
\setlength{\emergencystretch}{2em}
\setlist{nosep,leftmargin=1.5em,topsep=3pt}
\setcounter{secnumdepth}{0}
\newcommand{\source}[1]{{\small\color{muted}Источник: \texttt{handwritten-answers.pdf}, с.~#1.}\par}
\newcommand{\ednote}[1]{\par{\small\color{muted}\textit{Примечание к расшифровке.} #1}\par}
\newcommand{\question}[2]{\subsection{Вопрос #1. #2}}
\newcommand{\RSS}{\operatorname{RSS}}
\newcommand{\MSE}{\operatorname{MSE}}
\newcommand{\MAE}{\operatorname{MAE}}
\newcommand{\RMSE}{\operatorname{RMSE}}
\newcommand{\code}[1]{\texttt{#1}}
\newcolumntype{Y}{>{\raggedright\arraybackslash}X}

\begin{document}
\begin{center}
{\LARGE\sffamily\bfseries Machine Learning -- Semester 2}\par
{\large Набор рукописных конспектов: лекции 1--6}
\end{center}
\vspace{2pt}
{\small Расшифровка 19 фотографий. Сохранены темы и номера вопросов из записей;
это не полный комплект ответов к зачёту. Зачёркнутые черновики не перенесены.
Обозначения унифицированы; исправления явных описок отмечены отдельно.
Ссылки на источник указывают страницы PDF со снимками.}
\vspace{10pt}

\section{Лекция 1. Основы машинного обучения}
\source{1--2}

\question{7}{Разметка}
\textbf{Разметка} -- правильные ответы для объектов обучающей выборки.
Она необходима для обучения с учителем.

\question{50}{Этапы работы с моделью}
\begin{enumerate}
\item Постановка задачи и сбор данных.
\item Разведочный анализ: матрица корреляции, поиск дисбаланса и выбросов.
\item Предобработка: пропуски, выбросы, кодирование номинальных и порядковых
признаков, масштабирование, работа с датами, токенизация и обработка текста.
\item Разбиение на \code{train}, \code{validation}, \code{test}.
\item Обучение: выбор метода и подбор его параметров.
\item Подбор гиперпараметров: валидация или кросс-валидация.
\item Оценка качества на \code{test}: например, MSE, $R^2$, Accuracy, $F_1$.
\item Окончательная модель: внедрение и мониторинг, наблюдение за качеством
на новых данных, переобучение.
\end{enumerate}

\question{1}{Обучение с учителем и без учителя}
\textbf{С учителем:} обучающая выборка состоит из входов $x_i\in X$
и ответов $y_i\in Y$. Для классификации множество ответов конечно;
для регрессии обычно $Y=\mathbb R$.

\textbf{Без учителя:} обучающая выборка состоит только из входов $x_i\in X$,
без ответов. Алгоритм ищет закономерности в самих входных данных.
Пример -- кластеризация.

\question{3}{Типы признаков}
Пространство входов представлено как произведение областей значений признаков:
\[
 X=X_1\times X_2\times\cdots\times X_d,\qquad x_j\in X_j.
\]
\begin{tabularx}{\linewidth}{@{}lY@{}}
\toprule
Тип & Область значений\\
\midrule
Номинальный & Конечное неупорядоченное множество.\\
Порядковый & Конечное упорядоченное множество.\\
Количественный & Числовые значения; в записи $X_j=\mathbb R$.\\
Бинарный & $X_j=\{0,1\}$, два значения.\\
\bottomrule
\end{tabularx}

\question{4}{Классификация и регрессия}
\textbf{Классификация} -- задача обучения с учителем, в которой множество
ответов конечно: $|Y|<\infty$.

\textbf{Регрессия} -- задача обучения с учителем, в которой ответы числовые;
в записи $Y=\mathbb R$.

\clearpage
\section{Лекция 2. Pandas и табличные данные}
\source{3--4}

\question{48}{Series, DataFrame и индексация}
\textbf{Pandas} -- библиотека Python для анализа табличных данных.

\textbf{Series} -- одномерный массив значений с метками.
Пример из записи:
\begin{verbatim}
s = pd.Series([1, 2, 3, 4, 5], index=['a', 'b', 'c', 'd', 'e'])
s['a']       # 1
s.iloc[0]    # 1
\end{verbatim}

\textbf{DataFrame} -- двумерная таблица с помеченными строками и столбцами,
аналог таблицы Excel. Столбцы могут содержать числа, строки или даты;
каждый столбец -- Series.

\begin{center}
\begin{tabular}{c|rrr}
 & A & B & C\\\hline
a & 1 & 2 & 3\\
b & 4 & 5 & 6\\
c & 7 & 8 & 9\\
d & 10 & 11 & 12\\
e & 13 & 14 & 15
\end{tabular}
\end{center}

Формы обращения к таблице:
\begin{verbatim}
df['col']
df.loc[row_label, column_label]
df.iloc[row_position, column_position]
df[['col1', 'col2']]
\end{verbatim}
\code{.loc} обращается по меткам, \code{.iloc} -- по целочисленным позициям.
\ednote{Имена аргументов в схематических примерах заменены описательными,
чтобы различить метки и позиции.}

\question{49}{CSV и чтение данных}
\textbf{CSV} (comma-separated values) -- формат представления таблицы:
\begin{enumerate}
\item Каждая строка -- строка таблицы.
\item Значения столбцов разделены разделителем: например, запятой или
точкой с запятой. Десятичный разделитель может отличаться.
\item Первая строка обычно содержит названия столбцов.
\end{enumerate}
В записи приведён пример CSV с разделителем \code{;} и десятичной запятой:
\begin{verbatim}
pd.read_csv(name, sep=';', decimal=',')
\end{verbatim}

\clearpage
\section{Лекция 3. Построение и отбор признаков}
\source{5--8}

\question{22}{PCA}
\textbf{PCA} (Principal Component Analysis) -- метод сокращения размерности.
Он переводит объекты из многомерного пространства в пространство меньшей
размерности. В записях дано геометрическое описание: сумма квадратов расстояний
от объектов до их проекций минимальна; новые признаки -- линейные комбинации старых.

PCA не увеличивает размерность. В заметке упомянуты поворот и перестановка
координат как преобразования, не создающие новых измерений. Для увеличения
размерности нужны новые признаки, например полиномиальные произведения
или бинарное кодирование.

\question{27}{Мультиколлинеарность и корреляция}
\textbf{Мультиколлинеарность} -- сильная зависимость между двумя и более
признаками в выборке. Для её изучения используют матрицу корреляций.

Коэффициент корреляции Пирсона для двух числовых признаков:
\[
 r_{xz}=\frac{\sum_{i=1}^{n}(x_i-\bar x)(z_i-\bar z)}
 {\sqrt{\sum_{i=1}^{n}(x_i-\bar x)^2}
  \sqrt{\sum_{i=1}^{n}(z_i-\bar z)^2}},\qquad r_{xz}\in[-1,1].
\]
В записи: $1$ -- прямая связь, $-1$ -- обратная, $0$ -- связи нет.
\ednote{Значение $r=0$ означает отсутствие линейной корреляции,
а не обязательно отсутствие любой зависимости.}

Ограничения парной корреляции:
\begin{enumerate}
\item Рассматривает только два признака за раз.
\item Не обнаруживает нелинейные зависимости.
\item Применяется к числовым признакам в приведённой формуле.
\item Не выявляет причинность: связь $A$ с $B$ может объясняться общим фактором $C$.
\end{enumerate}

На полях отмечены возможная необратимость $X^\mathsf T X$ и регуляризация:
\[
 (X^\mathsf T X+\lambda I)^{-1}X^\mathsf T y.
\]

\question{32}{Число признаков и One-Hot Encoding}
Рост числа признаков увеличивает число параметров. Для той же выборки
это связано с риском переобучения.

\textbf{One-Hot Encoding} -- кодирование номинального категориального признака:
для каждого значения создаётся столбец, равный $1$ для этой категории
и $0$ в остальных случаях. Такой столбец задаёт разбиение выборки.
Если каждый объект принадлежит ровно одной категории, один из столбцов
можно убрать как однозначно определяемый остальными.

Преимущество -- не вводится ложный порядок. Недостаток -- много столбцов.
В сравнении с PCA такое кодирование может увеличивать размерность.

\textbf{Target encoding} (кодирование средним) -- замена категории средним
значением целевого признака для этой категории по выборке.

\question{33}{Порядковый признак и числовое кодирование}
\textbf{Порядковый признак} принимает упорядоченные значения.
В заметке объясняется недостаток One-Hot Encoding: модель не получает
порядок категорий непосредственно из кодов.

\textbf{Label / Ordinal Encoding} -- каждой категории присваивается числовой код
вроде $0,1,\ldots$.
\ednote{В рукописи Label Encoding и Ordinal Encoding объединены одной строкой.
Здесь сохранена эта группировка, без утверждения, что методы тождественны.}

\question{34}{Краткое повторение кодирований}
В рукописи перечислены порядковые и номинальные признаки,
Label Encoding и One-Hot Encoding; определения оставлены многоточиями.
Рядом указаны функции Pandas:
\begin{verbatim}
pd.factorize(...)
pd.get_dummies(...)
\end{verbatim}

\question{35}{Признаки из даты}
\begin{enumerate}
\item Число дней от начала наблюдений: помогает учитывать тренд.
\item Сезонные признаки: день недели, месяц, год, время суток.
\item Признаки особых дней.
\end{enumerate}
\ednote{Третий пункт уточнён автором рукописи.}

Внизу записи: \textbf{тренд} -- долговременное изменение целевой величины;
\textbf{сезонность} -- периодическое изменение целевой величины.

\question{36}{Бинирование}
\textbf{Бинирование} -- группировка количественного признака в категории
по диапазонам его значений. Пример -- разбиение возрастов порогами $18$ и $65$.

\question{37}{Пропуски и выбросы}
Способы обработки пропусков, перечисленные в записи:
\begin{enumerate}
\item Удаление, если пропусков мало.
\item Статистика: среднее, медиана; для категориального признака -- мода.
\item Предсказание пропущенного значения по другим признакам.
\item Специальная метка вроде \code{missing}, например для дерева решений,
и бинарный признак наличия пропуска.
\end{enumerate}
\[
 \bar x=\frac1n\sum_{i=1}^{n}x_i.
\]
Медиана -- середина упорядоченного ряда; при чётном числе элементов --
среднее двух центральных.

\textbf{Выбросы} анализируют по распределению каждого признака
или по всему вектору признаков. Нужно различать ошибочные значения
и верно измеренные, но нетипичные объекты.

\clearpage
\section{Лекция 4. Регрессия и переобучение}
\source{9--12}

\question{5}{Переобучение}
\textbf{Переобучение} -- ситуация, когда модель запоминает обучающую выборку
вместе с шумом и плохо обобщает на новые данные.

\textbf{Шум} -- случайная часть данных, которую нельзя предсказать по признакам:
\[
 y=f(x)+\varepsilon.
\]
Переобучение распознаётся по малой ошибке на \code{train}
и большой ошибке на \code{test}.

\question{6}{Почему нельзя оценивать качество только на train}
Ошибка на обучающей выборке обычно занижена и не показывает качество
на новых данных. Переобученная модель может выглядеть идеальной на \code{train}.
Оценивать итоговое качество нужно на данных, которых модель не видела.
Гиперпараметры подбираются по \code{validation} или кросс-валидацией.

\question{10}{Разбиение данных}
Данные разделяют на обучающую и тестовую выборки, чтобы оценить качество
на новых объектах. В записи приведены пропорции $70/30$ и $80/20$.

Слишком большой \code{train} оставляет мало данных для оценки;
слишком маленький -- мало данных для обучения.

\question{11}{Коэффициент детерминации}
$R^2$ показывает, какую долю разброса ответов объясняет модель
по сравнению с постоянным прогнозом, равным среднему:
\[
 R^2=1-\frac{\RSS}{\operatorname{TSS}}
 =1-\frac{\sum_{i=1}^{n}(y_i-\hat y_i)^2}
 {\sum_{i=1}^{n}(y_i-\bar y)^2}.
\]
\begin{itemize}
\item $R^2=1$ -- идеальное совпадение прогнозов и ответов.
\item $R^2=0$ -- качество как у прогноза средним.
\item $R^2<0$ -- хуже прогноза средним.
\end{itemize}

\question{28}{Потери, гиперпараметры и регуляризация}
Функция потерь сравнивает прогноз модели с правильным ответом
и выражает ошибку на одном примере:
\[
 L_1(y,\hat y)=|y-\hat y|,\qquad L_2(y,\hat y)=(y-\hat y)^2.
\]
\textbf{Гиперпараметры} -- настройки, задаваемые до обучения;
модель сама не подбирает их в процессе подгонки своих параметров.

Большие веса могут означать неустойчивость модели к вариациям входов.
Регуляризация добавляет штраф:
\begin{align*}
 \text{Lasso:}\quad &\RSS(w)+\lambda\sum_j|w_j|,\\
 \text{Ridge:}\quad &\RSS(w)+\lambda\sum_jw_j^2.
\end{align*}
Штраф уменьшает веса; Lasso может занулять часть коэффициентов.

\question{29}{Эластичная сеть}
\textbf{Elastic Net} использует оба штрафа:
\[
 \RSS(w)+\lambda_1\sum_j|w_j|+\lambda_2\sum_jw_j^2.
\]
Она сочетает отбор признаков от $L_1$ и устойчивость к корреляциям от $L_2$.
В записи отмечено: при коррелирующих признаках Lasso может нестабильно
выбирать один из них, а Ridge распределяет веса между признаками.

\question{30}{Большие коэффициенты линейной регрессии}
Причины для проверки:
\begin{enumerate}
\item Переобучение: модель может подгоняться под шум.
\item Мультиколлинеарность: у коррелирующих признаков возможны большие веса
разных знаков.
\item Разные масштабы признаков: величину коэффициента нельзя интерпретировать
без учёта единиц измерения входа.
\end{enumerate}

\question{31}{Метод наименьших квадратов}
МНК находит параметры линейной регрессии минимизацией суммы квадратов ошибок:
\[
 w^*=\arg\min_w\RSS(w),\qquad
 \RSS(w)=\sum_{i=1}^{n}(y_i-w^\mathsf T x_i)^2.
\]
Приведённое в рукописи решение:
\[
 w^*=(X^\mathsf T X)^{-1}X^\mathsf T y.
\]
\ednote{В строке с RSS на с.~12 исходного PDF пропущен квадрат;
он восстановлен по названию метода и остальным записям.
Формула через обратную матрицу предполагает её существование.}

Выбросы вносят большой вклад в RSS. Гиперплоскость может притягиваться
к выбросу и хуже описывать остальную зависимость.

\clearpage
\section{Лекция 5. Ближайшие соседи и метрики}
\source{13--15}

\question{9}{Метод ближайших соседей}
$k$-NN -- метод классификации и регрессии. Для нового объекта выбираются
$k$ ближайших объектов обучающей выборки по заданной мере близости.
В классификации используется наиболее частый класс среди соседей,
в регрессии -- усреднение ответов.

Метод непараметрический: хранит обучающую выборку. В варианте
с голосованием по меткам ему нужна размеченная выборка, поэтому это обучение
с учителем.

Расстояния из записи:
\begin{align*}
 \rho_2(x,z)&=\sqrt{\sum_j(x_j-z_j)^2}, &&\text{евклидово},\\
 \rho_{\cos}(x,z)&=1-\frac{x^\mathsf Tz}{\|x\|\,\|z\|}, &&\text{косинусное},\\
 \rho_1(x,z)&=\sum_j|x_j-z_j|, &&\text{манхэттенское}.
\end{align*}
\ednote{В последней формуле на фото не видны модули;
они восстановлены для манхэттенского расстояния.}

\question{14}{Метрики регрессии и классификации}
\begin{align*}
 \operatorname{SAE}&=\sum_{i=1}^{n}|y_i-\hat y_i|,\\
 \MAE&=\frac1n\sum_{i=1}^{n}|y_i-\hat y_i|,\\
 \MSE&=\frac1n\sum_{i=1}^{n}(y_i-\hat y_i)^2,\\
 \RMSE&=\sqrt{\MSE}.
\end{align*}

Обозначения матрицы ошибок: TP и TN -- правильные положительные
и отрицательные ответы; FP -- ложноположительные, FN -- ложноотрицательные.

\begin{center}
\begin{tabular}{lcc}
\toprule
 & Прогноз $1$ & Прогноз $0$\\
\midrule
Истина $1$ & TP & FN\\
Истина $0$ & FP & TN\\
\bottomrule
\end{tabular}
\end{center}
\ednote{Оси матрицы подписаны явно, чтобы не воспроизводить
неоднозначность схематического расположения ячеек на фотографии.}

\begin{align*}
 \operatorname{Accuracy}&=\frac{TP+TN}{TP+TN+FP+FN},\\
 \operatorname{Precision}&=\frac{TP}{TP+FP},\\
 \operatorname{Recall}&=\frac{TP}{TP+FN},\\
 F_1&=2\frac{\operatorname{Precision}\operatorname{Recall}}
 {\operatorname{Precision}+\operatorname{Recall}}.
\end{align*}
$F_1$ -- гармоническое среднее Precision и Recall.
В записи отмечена идея баланса между этими метриками.

\question{15}{Precision и Recall}
\textbf{Precision} -- доля правильных положительных прогнозов среди всех
положительных прогнозов. Важна, когда нужно быть уверенным в положительном
решении; пример в записи -- спам-фильтр. Уменьшают число FP.

\textbf{Recall} -- доля найденных положительных объектов среди всех
действительно положительных. Важна, когда нельзя пропустить объект;
пример -- диагностика болезни. Уменьшают число FN.

В записи FP названы ошибками I рода, FN -- ошибками II рода.
Метрики могут конкурировать; $F_1$ отражает их баланс.

\question{16}{Дисбаланс классов и Accuracy}
Прогноз всегда самого большого класса может дать высокую Accuracy,
хотя модель не решает задачу обнаружения другого класса.

\clearpage
\section{Лекция 6. Деревья решений и ансамбли}
\source{16--19}

\question{23}{Построение дерева решений}
\textbf{Дерево решений} решает задачи классификации и регрессии:
данные последовательно делятся на группы по значениям признаков.
Прогноз для нового объекта -- путь от корня к листу по правилам разбиения;
ответ определяется листом. В записи рассматриваются бинарные одномерные деревья.

Построение одного дерева:
\begin{enumerate}
\item Выбрать лучшее разбиение объектов текущего узла.
Для числового признака проверяются пары «признак, порог».
Если $x_1<x_2<\cdots<x_n$ -- различные упорядоченные значения, кандидаты:
\[
 c_i=\frac{x_i+x_{i+1}}2,\qquad i=1,\ldots,n-1.
\]
Для категориального признака рассматриваются разбиения категорий на две группы;
в записи указано $2^{k-1}-1$ вариантов для $k$ категорий.
\item Разделить объекты на дочерние узлы.
\item Рекурсивно повторить выбор разбиения в каждом узле.
\item Остановиться по условию остановки.
\item Присвоить листьям прогнозы.
\end{enumerate}

Качество разбиения -- взвешенная неоднородность дочерних узлов:
\[
 Q=\frac{n_L}{n}C(L)+\frac{n_R}{n}C(R),\qquad n=n_L+n_R.
\]
Для классификации приведены:
\[
 H=-\sum_jp_j\log_2p_j,\qquad
 \operatorname{Gini}=1-\sum_jp_j^2.
\]
Энтропия описывает неопределённость класса случайного объекта.
Индекс Джини связан с вероятностью несовпадения классов при случайном выборе
классов по их долям в узле.

Для регрессии в записи используются MSE и MAE:
\[
 \MSE=\frac1n\sum_i(y_i-\bar y)^2,\qquad
 \MAE=\frac1n\sum_i|y_i-\widetilde y|,
\]
где $\bar y$ -- среднее, $\widetilde y$ -- медиана.
\ednote{В строке MAE на с.~16 исходного PDF ошибочно стоит квадрат;
он заменён модулем, как и в формуле MAE на с.~14 исходного PDF.}

Прирост информации и Gain Ratio:
\begin{align*}
 \operatorname{Gain}&=H(\text{узел})-Q,\\
 \operatorname{GainRatio}&=\frac{\operatorname{Gain}}{\operatorname{SplitInfo}},\\
 \operatorname{SplitInfo}&=-\sum_k\frac{n_k}{n}\log_2\frac{n_k}{n}.
\end{align*}
SplitInfo -- энтропия самого разбиения. Для регрессии критерий можно
представить как уменьшение ошибки:
\[
 \Delta=\MSE(\text{узел})-Q.
\]
Выбирается разбиение, сильнее уменьшающее разброс ответов в группах.

Условия остановки из записи: чистая группа объектов, максимальная глубина,
минимальное число объектов и отсутствие улучшения качества.
Прогноз в листе: наиболее частый класс для классификации, среднее для регрессии.

\question{24}{Масштаб признаков}
Одиночные одномерные деревья нечувствительны к масштабу количественного признака:
сравнение $x_j\le c$ зависит от порядка значений, а не от их величины.
Комбинации признаков в рассматриваемом варианте дерева не используются.

\question{25}{Ансамбли}
\textbf{Ансамбль} -- сочетание нескольких алгоритмов для получения более
эффективной модели. В записи перечислены бэггинг, бустинг и стекинг;
стекинг отдельно не раскрыт.

\textbf{Бэггинг} (bootstrap aggregating): модели обучаются параллельно
и независимо на bootstrap-выборках -- случайных выборках с возвращением.
Прогнозы объединяются усреднением или голосованием.
Ошибки разных моделей частично компенсируются, уменьшая вариативность.

\textbf{Бустинг}: модели строятся последовательно; каждая новая модель
исправляет ошибки предыдущих, усиливая итоговый ансамбль.

\textbf{Случайный лес}: ансамбль деревьев по принципу бэггинга.
При выборе разбиения в узле рассматривается случайное подмножество признаков;
в записи приведён ориентир $m\approx\sqrt d$.

\textbf{Градиентный бустинг}: каждая новая модель обучается предсказывать
остатки текущего ансамбля при квадратичной функции потерь, затем добавляется
с небольшим шагом. Итоговый прогноз -- сумма начального прогноза и вкладов моделей.
\ednote{В рукописи написано «остатки» для градиентного бустинга вообще.
Оговорка о квадратичной функции потерь добавлена, чтобы не приравнивать
остатки к антиградиенту для любой функции потерь.}

\begin{tabularx}{\linewidth}{@{}lYY@{}}
\toprule
 & Случайный лес & Градиентный бустинг\\
\midrule
Базовая модель & Дерево решений & Любая базовая модель\\
Принцип & Бэггинг & Бустинг\\
Объединение & Усреднение или голосование & Взвешенная сумма вкладов\\
\bottomrule
\end{tabularx}

\end{document}
